\begin{abstract}
A capability regime in which agents continuously discover and
escalate new capabilities faster than any generation can experience
their risks through ordinary operation forces a sharper question than
the classical risk-management literature addresses: how do we preserve
safety without requiring the incalculable cost-of-life that produced
the historical markers of prior eras---the tsunami stones at Fudai
that warned, in stone and across centuries, against construction
below the inundation line?  The historical-markers mechanism requires
a first catastrophe to write the markers.  In a capability landscape
whose technological cadence outpaces any community's lived
experience, the first catastrophe is the last opportunity to learn
the lesson, and the markers that would have warned the next
generation cannot be relied on to exist.  Risks with realisation
timelines far exceeding the practical observational window produce a
structural pathology in the Goal-Frontier Maximization risk
framework: a restriction justified by such a risk cannot be revised
through observation because the evidence that would revise it cannot
arrive within the decision cycle, and the alternative mechanism
---waiting for the catastrophe that would make the evidence
available---is the failure mode we are trying to avoid.  We call this
the \emph{Wamura problem}, after the tsunami-preparedness case at
Fudai; the historical case is motivating rather than directly soluble
by our construction, because Fudai had ignored historical evidence
available (centuries-old tsunami stones) rather than lacking it.  The
paper addresses the harder complementary case, in which no such
historical record exists or can be expected to exist on the
technology's cadence.  We formalize the structural
asymmetry, prove that max-aggregation with neutral-prior initialization
produces monotone overcaution absent external intervention, and
introduce \emph{controlled relaxation protocols}: bounded, monitored
exercise windows that generate genuine Bayesian evidence about risk
claims while provably bounding maximum damage.  Two central results
follow.  First, a scope-construction theorem bounds the damage from a
test by a user-set parameter $\damagebound$ even when cross-boundary
cooperative capabilities could propagate the test's effects.  Second,
under a simple binary hypothesis framework with fixed likelihoods, a
convergence theorem shows that repeated controlled relaxation drives
risk estimates almost surely to truth at exponential rate, with the
rate given by the Kullback--Leibler divergence between the observation
distributions.  We further extend the protocol to
\emph{keystone-refinement tests}, which handle the common real-world
case where a restriction is not obsolete but misconfigured---the
underlying risk pathway still operates, but technological progress has
made narrower keystones available that would block it without blocking
adjacent benign pathways.  The protocol closes the most important open
problem named in~\citep{lasser2026horizon} and resolves the
degenerate-phase-boundary case of~\citep[Proposition~4]{lasser2026phase}.
\end{abstract}
